% !PRINTED 318
\weberpage{318}
%<<CONT
auch $T$ intransitiv. Daraus folgt, wie oben gezeigt, dass $TQ$
imprimitiv ist.

Als Gruppe für primitive metacyklische Gleichungen 8\textsuperscript{ten}
Grades bleiben also nur die beiden Typen
\[
T_7 Q, \quad T_{21} Q
\]
übrig. Im ersten Falle sind die in der Formel [\S.\,82, (3)] vorkommenden
Grössen $v$ die Wurzeln einer cyklischen Gleichung 7\textsuperscript{ten}
Grades mit der Gruppe $(z, z+b)$, im zweiten einer metacyklischen mit der
Gruppe $(z, az+b)$ (Bd. I, \S.\,180), worin $z, a, b$ nach dem Modul 7
genommen sind und $a$ die Werthe 1, 2, 4 durchläuft.

Um eine Zuordnung der $v_i$, wie sie in der Formel (3), \S.\,82 vorkommen,
zu den $v_z$, worin $z$ der in den Gruppen $(z, z+b)$ oder $(z, az+b)$
vorkommende Index ist, zu erhalten, können wir von einer beliebigen
Substitution siebenter Ordnung ausgehen, etwa von
\[
\tau = \begin{pmatrix} 1, & 0, & 1 \\ 1, & 0, & 0 \\ 0, & 1, & 0 \end{pmatrix}
\quad [\S.\,79, (2)].
\]

Nehmen wir dann $v_{1,0,0}$ für $v_0$, so findet man durch wiederholte
Anwendung von $\tau$ auf $1, 0, 0$ folgende Anordnung:
\[
(1,0,0)=0, \quad (1,1,0)=1, \quad (1,1,1)=2, \quad (0,1,1)=3, \quad (1,0,1)=4,
\]
\[
(0,1,0)=5, \quad (0,0,1)=6,
\]
und nach \S.\,82, (2) hat man daher
\[
v_1, v_2, v_3, v_4, v_5, v_6, v_7
\]
durch
\[
v_0, v_5, v_6, v_3, v_4, v_1, v_2
\]
zu ersetzen. Dadurch ergiebt sich aus \S.\,82, (3):
\[
8X = A \sum_{0,\,6}^{z} \sqrt{v_z}\,\sqrt{v_{z+1}}\,\sqrt{v_{z+2}}\,\sqrt{v_{z+4}} + B,
\]
worin $A, B$ rationale Grössen, und $v_z$ die Wurzeln einer cyklischen oder
einer metacyklischen Gleichung von der Gruppe $(z, az+b)$ vom Grade 7 sind.

Umgekehrt ist es auch nicht schwer, nachzuweisen, dass die acht durch
Wechsel der Vorzeichen hieraus abgeleiteten Grössen immer die Wurzeln
einer primitiven metacyklischen Gleichung
%CONT>>
